\documentclass{article}
\usepackage{amsmath, amssymb, geometry, xcolor}
\geometry{a4paper, margin=1in}

\title{Homework 3.2: Properties of Determinants}
\author{Tashfeen Omran}
\date{September 2026}

\begin{document}

\maketitle
\tableofcontents
\newpage

\section{Part 1: Comprehensive Introduction, Context, and Prerequisites}

\subsection{Core Concepts}
The determinant is a scalar value that can be computed from the elements of a square matrix and encodes certain core properties of the linear transformation described by the matrix. In Section 3.2, we transition from the mechanical computation of determinants (via cofactor expansion) to understanding their algebraic properties. The most crucial concept is how elementary row operations affect the determinant. By understanding these properties, we can vastly simplify the computation of determinants for large matrices by reducing them to upper triangular (echelon) form, where the determinant is simply the product of the diagonal entries.

\subsection{Intuition and Derivation}
Geometrically, the absolute value of the determinant of an $n \times n$ matrix represents the volume of the $n$-dimensional parallelepiped spanned by its column vectors. The sign of the determinant indicates the orientation of this volume (e.g., whether the transformation preserves or reverses the "handedness" of the space). 
When we perform row operations, we are altering this geometry in predictable ways:
\begin{itemize}
    \item \textbf{Row Replacement ($R_i \leftarrow R_i + cR_j$):} This corresponds to a shear transformation. Shearing shifts the shape but does not change its base or height, thus the volume (and the determinant) remains unchanged.
    \item \textbf{Row Interchange ($R_i \leftrightarrow R_j$):} This reflects the space across a plane, reversing the orientation. Thus, the determinant is multiplied by $-1$.
    \item \textbf{Scalar Multiplication ($R_i \leftarrow cR_i$):} This stretches the volume in one dimension by a factor of $c$. Thus, the determinant is multiplied by $c$.
\end{itemize}

\subsection{Historical Context and Motivation}
The concept of the determinant actually predates the modern definition of a matrix. In the 17th and 18th centuries, mathematicians like Gottfried Leibniz and Gabriel Cramer were searching for general formulas to solve systems of linear equations. They discovered that the existence of a unique solution was "determined" by a specific polynomial combination of the system's coefficients—hence the name "determinant." If this value was non-zero, a unique solution existed. This historical motivation directly links to Theorem 4 in this section: A matrix is invertible if and only if its determinant is non-zero.

\subsection{Key Formulas}
\begin{itemize}
    \item \textbf{Row Operations:} 
    \begin{itemize}
        \item Replacement: $\det(A_{\text{replace}}) = \det(A)$
        \item Interchange: $\det(A_{\text{swap}}) = -\det(A)$
        \item Scaling: $\det(A_{\text{scale}}) = c \det(A)$
    \end{itemize}
    \item \textbf{Invertibility:} $A$ is invertible $\iff \det(A) \neq 0$.
    \item \textbf{Transpose Property:} $\det(A^T) = \det(A)$.
    \item \textbf{Multiplicative Property:} $\det(AB) = \det(A)\det(B)$.
    \item \textbf{Inverse Property:} $\det(A^{-1}) = \frac{1}{\det(A)}$.
    \item \textbf{Scalar Multiplication of a Matrix:} $\det(cA) = c^n \det(A)$ for an $n \times n$ matrix.
\end{itemize}

\subsection{Prerequisites}
To master this section, you must be proficient in:
\begin{itemize}
    \item \textbf{Gaussian Elimination:} Performing elementary row operations to reduce a matrix to echelon form.
    \item \textbf{Cofactor Expansion:} Computing determinants using the recursive definition from Section 3.1.
    \item \textbf{Matrix Algebra:} Understanding matrix multiplication, inverses, and transposes.
\end{itemize}

\newpage
\section{Part 2: Detailed Homework Solutions}
\textit{Note: A strategic selection of problems has been chosen to cover every distinct problem archetype found in the text, ensuring full conceptual coverage without redundant repetition.}

\subsection{Problem 5: Determinant by Row Reduction}
\textbf{Question:} Find the determinant by row reduction to echelon form:
\[ \begin{vmatrix} 1 & 5 & -4 \\ -1 & -4 & 5 \\ -2 & -8 & 7 \end{vmatrix} \]

\textbf{Solution:}
We will use row replacement operations to create zeros below the first pivot (the $1$ in the top-left corner). Row replacements do not change the value of the determinant.
\begin{align*}
    \begin{vmatrix} 1 & 5 & -4 \\ -1 & -4 & 5 \\ -2 & -8 & 7 \end{vmatrix} 
    &= \begin{vmatrix} 1 & 5 & -4 \\ 0 & 1 & 1 \\ -2 & -8 & 7 \end{vmatrix} \quad \text{(using } R_2 \leftarrow R_2 + R_1 \text{)} \\
    &= \begin{vmatrix} 1 & 5 & -4 \\ 0 & 1 & 1 \\ 0 & 2 & -1 \end{vmatrix} \quad \text{(using } R_3 \leftarrow R_3 + 2R_1 \text{)}
\end{align*}
Now, we create a zero below the second pivot (the $1$ in the middle):
\begin{align*}
    &= \begin{vmatrix} 1 & 5 & -4 \\ 0 & 1 & 1 \\ 0 & 0 & -3 \end{vmatrix} \quad \text{(using } R_3 \leftarrow R_3 - 2R_2 \text{)}
\end{align*}
The matrix is now in upper triangular (echelon) form. The determinant of an upper triangular matrix is the product of its diagonal entries.
\[ \det(A) = (1)(1)(-3) = -3 \]
\textbf{Final Answer:} $-3$

\subsection{Problem 11: Combining Row Reduction and Cofactor Expansion}
\textbf{Question:} Combine the methods of row reduction and cofactor expansion to compute the determinant:
\[ \begin{vmatrix} 3 & 4 & -3 & -1 \\ 3 & 0 & 1 & -3 \\ -6 & 0 & -4 & 3 \\ 6 & 8 & -4 & -1 \end{vmatrix} \]

\textbf{Solution:}
Notice that the second column already has two zeros. We can create a third zero in the second column to make cofactor expansion trivial. We will use the $4$ in row 1 to eliminate the $8$ in row 4.
\begin{align*}
    \begin{vmatrix} 3 & 4 & -3 & -1 \\ 3 & 0 & 1 & -3 \\ -6 & 0 & -4 & 3 \\ 6 & 8 & -4 & -1 \end{vmatrix} 
    &= \begin{vmatrix} 3 & 4 & -3 & -1 \\ 3 & 0 & 1 & -3 \\ -6 & 0 & -4 & 3 \\ 0 & 0 & 2 & 1 \end{vmatrix} \quad \text{(using } R_4 \leftarrow R_4 - 2R_1 \text{)}
\end{align*}
Now, expand along the second column. The only non-zero term is the $4$ at position $(1,2)$. The sign for this position is $(-1)^{1+2} = -1$.
\[ = 4(-1) \begin{vmatrix} 3 & 1 & -3 \\ -6 & -4 & 3 \\ 0 & 2 & 1 \end{vmatrix} = -4 \begin{vmatrix} 3 & 1 & -3 \\ -6 & -4 & 3 \\ 0 & 2 & 1 \end{vmatrix} \]
Next, we can use row reduction on this $3 \times 3$ matrix to create a zero in the first column.
\begin{align*}
    -4 \begin{vmatrix} 3 & 1 & -3 \\ -6 & -4 & 3 \\ 0 & 2 & 1 \end{vmatrix} 
    &= -4 \begin{vmatrix} 3 & 1 & -3 \\ 0 & -2 & -3 \\ 0 & 2 & 1 \end{vmatrix} \quad \text{(using } R_2 \leftarrow R_2 + 2R_1 \text{)}
\end{align*}
Expand along the first column:
\[ = -4 \left( 3 \begin{vmatrix} -2 & -3 \\ 2 & 1 \end{vmatrix} \right) = -12 \begin{vmatrix} -2 & -3 \\ 2 & 1 \end{vmatrix} \]
Evaluate the $2 \times 2$ determinant:
\[ = -12 ((-2)(1) - (-3)(2)) = -12 (-2 + 6) = -12(4) = -48 \]
\textbf{Final Answer:} $-48$

\subsection{Problem 15 \& 19: Abstract Determinant Properties}
\textbf{Question 15:} Given $\begin{vmatrix} a & b & c \\ d & e & f \\ g & h & i \end{vmatrix}$, find $\begin{vmatrix} a & b & c \\ d & e & f \\ 3g & 3h & 3i \end{vmatrix}$.

\textbf{Solution 15:}
By Theorem 3(c), if one row of a matrix is multiplied by $k$, the determinant is multiplied by $k$. Here, the third row is multiplied by $3$.
\[ \begin{vmatrix} a & b & c \\ d & e & f \\ 3g & 3h & 3i \end{vmatrix} = 3 \begin{vmatrix} a & b & c \\ d & e & f \\ g & h & i \end{vmatrix} \]

\textbf{Question 19:} Given $\begin{vmatrix} a & b & c \\ d & e & f \\ g & h & i \end{vmatrix}$, find $\begin{vmatrix} a & b & c \\ 2d+a & 2e+b & 2f+c \\ g & h & i \end{vmatrix}$.

\textbf{Solution 19:}
First, we can apply a row replacement operation in reverse. Subtracting row 1 from row 2 does not change the determinant:
\[ \begin{vmatrix} a & b & c \\ 2d+a & 2e+b & 2f+c \\ g & h & i \end{vmatrix} = \begin{vmatrix} a & b & c \\ 2d & 2e & 2f \\ g & h & i \end{vmatrix} \quad \text{(using } R_2 \leftarrow R_2 - R_1 \text{)} \]
Next, factor out the scalar $2$ from the second row:
\[ = 2 \begin{vmatrix} a & b & c \\ d & e & f \\ g & h & i \end{vmatrix} \]

\subsection{Problem 21: Invertibility via Determinants}
\textbf{Question:} Use determinants to find out if the matrix is invertible:
\[ A = \begin{bmatrix} 1 & 3 & 6 \\ 2 & 4 & 7 \\ 0 & 5 & 8 \end{bmatrix} \]

\textbf{Solution:}
By Theorem 4, $A$ is invertible if and only if $\det(A) \neq 0$. We compute the determinant by expanding along the first column:
\begin{align*}
    \det(A) &= 1 \begin{vmatrix} 4 & 7 \\ 5 & 8 \end{vmatrix} - 2 \begin{vmatrix} 3 & 6 \\ 5 & 8 \end{vmatrix} + 0 \\
    &= 1(32 - 35) - 2(24 - 30) \\
    &= 1(-3) - 2(-6) \\
    &= -3 + 12 = 9
\end{align*}
Since $\det(A) = 9 \neq 0$, the matrix is invertible.
\textbf{Final Answer:} The matrix is invertible.

\subsection{Problem 24: Linear Independence via Determinants}
\textbf{Question:} Use determinants to decide if the set of vectors is linearly independent:
\[ \begin{bmatrix} 4 \\ 6 \\ 2 \end{bmatrix}, \begin{bmatrix} -6 \\ 0 \\ 6 \end{bmatrix}, \begin{bmatrix} -3 \\ -5 \\ -2 \end{bmatrix} \]

\textbf{Solution:}
A set of $n$ vectors in $\mathbb{R}^n$ is linearly independent if and only if the matrix formed by these vectors as columns is invertible, which means its determinant must be non-zero.
Let $A = \begin{bmatrix} 4 & -6 & -3 \\ 6 & 0 & -5 \\ 2 & 6 & -2 \end{bmatrix}$.
We expand along the second row to take advantage of the zero:
\begin{align*}
    \det(A) &= -6 \begin{vmatrix} -6 & -3 \\ 6 & -2 \end{vmatrix} + 0 - (-5) \begin{vmatrix} 4 & -6 \\ 2 & 6 \end{vmatrix} \\
    &= -6(12 - (-18)) + 5(24 - (-12)) \\
    &= -6(30) + 5(36) \\
    &= -180 + 180 = 0
\end{align*}
Since $\det(A) = 0$, the matrix is not invertible, meaning its columns are linearly dependent.
\textbf{Final Answer:} The vectors are linearly dependent.

\subsection{Problem 37 \& 40: Theoretical Proofs}
\textbf{Question 37:} Show that if $A$ is invertible, then $\det(A^{-1}) = \frac{1}{\det(A)}$.

\textbf{Solution 37:}
If $A$ is invertible, then by definition, $A A^{-1} = I$, where $I$ is the identity matrix.
Taking the determinant of both sides yields:
\[ \det(A A^{-1}) = \det(I) \]
By the Multiplicative Property (Theorem 6), $\det(AB) = \det(A)\det(B)$. Also, the determinant of the identity matrix is $1$.
\[ \det(A) \det(A^{-1}) = 1 \]
Since $A$ is invertible, Theorem 4 guarantees that $\det(A) \neq 0$. Therefore, we can safely divide both sides by $\det(A)$:
\[ \det(A^{-1}) = \frac{1}{\det(A)} \]

\textbf{Question 40:} Let $A$ and $P$ be square matrices, with $P$ invertible. Show that $\det(PAP^{-1}) = \det(A)$.

\textbf{Solution 40:}
We apply the Multiplicative Property of determinants ($\det(XY) = \det(X)\det(Y)$) repeatedly:
\begin{align*}
    \det(PAP^{-1}) &= \det(P) \det(A) \det(P^{-1})
\end{align*}
Using the result from Problem 37, we substitute $\det(P^{-1}) = \frac{1}{\det(P)}$:
\begin{align*}
    &= \det(P) \det(A) \frac{1}{\det(P)}
\end{align*}
Since determinants are scalar values, multiplication is commutative. The $\det(P)$ and $\frac{1}{\det(P)}$ cancel out:
\begin{align*}
    &= \det(A) \left( \det(P) \frac{1}{\det(P)} \right) = \det(A) \cdot 1 = \det(A)
\end{align*}

\subsection{Problem 45: Multiplicative Properties in Practice}
\textbf{Question:} Let $A$ and $B$ be $3 \times 3$ matrices, with $\det(A) = -2$ and $\det(B) = 3$. Compute:
a. $\det(AB)$ \quad b. $\det(5A)$ \quad c. $\det(B^T)$ \quad d. $\det(A^{-1})$ \quad e. $\det(A^3)$

\textbf{Solution:}
\begin{itemize}
    \item \textbf{a.} By Theorem 6, $\det(AB) = \det(A)\det(B) = (-2)(3) = -6$.
    \item \textbf{b.} When multiplying an $n \times n$ matrix by a scalar $c$, the determinant is multiplied by $c^n$. Since $A$ is $3 \times 3$, $\det(5A) = 5^3 \det(A) = 125(-2) = -250$.
    \item \textbf{c.} By Theorem 5, $\det(B^T) = \det(B) = 3$.
    \item \textbf{d.} As proven in Ex 37, $\det(A^{-1}) = \frac{1}{\det(A)} = -\frac{1}{2}$.
    \item \textbf{e.} $\det(A^3) = \det(A \cdot A \cdot A) = \det(A)\det(A)\det(A) = (\det(A))^3 = (-2)^3 = -8$.
\end{itemize}

\newpage
\section{Part 3: In-Depth Analysis of Problems and Techniques}

\subsection{A) Problem Types and General Approach}
\begin{itemize}
    \item \textbf{Type 1: Determinant via Row Reduction (Problems 5, 11):} 
    \textit{Strategy:} Use row replacement operations ($R_i \leftarrow R_i + cR_j$) to create zeros below the main diagonal. Because this specific operation does not change the determinant, you can reduce the matrix to upper triangular form and simply multiply the diagonal entries. If you must use a row swap, remember to multiply the final result by $-1$.
    \item \textbf{Type 2: Abstract Row Operations (Problems 15, 19):}
    \textit{Strategy:} Work backward from the modified matrix to the original matrix. Identify which elementary row operations were applied and adjust the determinant accordingly (e.g., factoring out a scalar from a single row).
    \item \textbf{Type 3: Invertibility and Linear Independence (Problems 21, 24):}
    \textit{Strategy:} This is an application of the Invertible Matrix Theorem. Compute the determinant using any method. If $\det(A) = 0$, the matrix is singular (not invertible) and its columns are linearly dependent. If $\det(A) \neq 0$, it is invertible and its columns are independent.
    \item \textbf{Type 4: Algebraic Properties of Determinants (Problems 37, 40, 45):}
    \textit{Strategy:} Apply the core theorems: $\det(AB) = \det(A)\det(B)$, $\det(A^T) = \det(A)$, and $\det(cA) = c^n \det(A)$. Treat the determinants as standard algebraic scalars once the properties are applied.
\end{itemize}

\subsection{B) Key Algebraic and Calculus Manipulations}
\begin{itemize}
    \item \textbf{Hybrid Evaluation (The "Best of Both Worlds" Trick):} Seen in Problem 11. Instead of fully row-reducing a $4 \times 4$ matrix (which can introduce messy fractions), we used just one row operation to create a column with mostly zeros, and then immediately used cofactor expansion. This minimizes arithmetic errors.
    \item \textbf{Factoring Scalars from a Single Row:} Seen in Problem 19. Unlike standard matrix scalar multiplication where $cA$ multiplies every entry by $c$, the determinant property allows you to factor a scalar out of a \textit{single row or column}. $\det(\text{row } i \times c) = c \det(A)$.
    \item \textbf{The Power of $n$ in Scalar Multiplication:} Seen in Problem 45b. A very common mistake is writing $\det(5A) = 5 \det(A)$. The correct manipulation requires recognizing that $5A$ multiplies \textit{every} row by $5$. For a $3 \times 3$ matrix, this means factoring out a $5$ three times, yielding $5^3 \det(A)$.
\end{itemize}

\newpage
\section{Part 4: "Cheatsheet" and Tips for Success}

\subsection{Summary of Important Formulas}
\begin{align*}
    \det(AB) &= \det(A)\det(B) \\
    \det(A^T) &= \det(A) \\
    \det(A^{-1}) &= \frac{1}{\det(A)} \quad \text{(if } A \text{ is invertible)} \\
    \det(cA) &= c^n \det(A) \quad \text{(where } A \text{ is an } n \times n \text{ matrix)} \\
    \det(I_n) &= 1
\end{align*}

\subsection{Cheats, Tricks, and Shortcuts}
\begin{itemize}
    \item \textbf{The Triangular Shortcut:} If a matrix is upper triangular, lower triangular, or diagonal, stop calculating! The determinant is just the product of the main diagonal entries.
    \item \textbf{The Zero Row/Column Cheat:} If a matrix has a row or column of all zeros, its determinant is $0$.
    \item \textbf{The Proportionality Cheat:} If two rows (or two columns) are identical, or if one is a scalar multiple of the other, the determinant is $0$. (You can prove this by doing one row replacement to create a row of zeros).
    \item \textbf{Avoid Fractions:} When row reducing to find a determinant, you do \textit{not} need to create leading $1$s. Creating leading $1$s often requires dividing a row by a scalar, which introduces fractions and changes the determinant. Just use the existing numbers to eliminate entries below them.
\end{itemize}

\subsection{Common Pitfalls and Traps}
\begin{itemize}
    \item \textbf{The Linearity Trap:} $\det(A + B) \neq \det(A) + \det(B)$. This is a massive misconception. Determinants are \textit{not} linear over matrix addition.
    \item \textbf{Forgetting the Sign Flip:} If you swap two rows to make row reduction easier, you \textit{must} multiply your final determinant by $-1$ for every swap you made.
    \item \textbf{Scalar Multiplication Error:} Forgetting the exponent $n$ when calculating $\det(cA)$.
\end{itemize}

\newpage
\section{Part 5: Conceptual Synthesis and The "Big Picture"}

\subsection{A) Thematic Connections}
The central theme of this topic is \textbf{invariance and transformation}. In mathematics, we frequently look for properties that remain unchanged (invariant) when we manipulate an object. Here, we discover that the "volume" (determinant) of a matrix is invariant under shear transformations (row replacements). This mirrors concepts in calculus, such as how the area under a curve remains invariant under certain coordinate shifts, or how energy is conserved in physics despite changes in state.

\subsection{B) Forward and Backward Links}
\begin{itemize}
    \item \textbf{Backward Links:} This section builds directly on Gaussian elimination (Chapter 1). We are using the exact same mechanical algorithms, but instead of using them to find a solution vector $x$, we are tracking how the algorithms affect the geometric volume of the matrix.
    \item \textbf{Forward Links:} 
    \begin{itemize}
        \item \textbf{Eigenvalues (Chapter 5):} The most important application of determinants is finding eigenvalues. We will solve the characteristic equation $\det(A - \lambda I) = 0$. Without the properties learned in this section, computing this polynomial for large matrices would be impossible.
        \item \textbf{Multivariable Calculus:} The determinant forms the core of the \textit{Jacobian matrix}, which is required for changing variables in multiple integrals (the multivariable equivalent of $u$-substitution). It scales the infinitesimal volume element $dx \, dy \, dz$.
    \end{itemize}
\end{itemize}

\newpage
\section{Part 6: Abstract Real-World Modeling (AI, ML, \& Quant Finance)}

\subsection{A) Concrete Scenarios}
\begin{enumerate}
    \item \textbf{Machine Learning (Gaussian Mixture Models):} In unsupervised learning, clustering data often relies on Multivariate Normal Distributions. The determinant of the covariance matrix $\Sigma$ appears in the denominator of the probability density function. It normalizes the distribution by measuring the $n$-dimensional spread (volume) of the data.
    \item \textbf{Information Theory (Differential Entropy):} The entropy of a continuous multivariate Gaussian random variable is directly proportional to $\ln(\det(\Sigma))$. A larger determinant means the data is more spread out, equating to higher uncertainty or entropy.
    \item \textbf{Quantitative Finance (Copula Models):} When pricing collateralized debt obligations (CDOs) or managing portfolio risk, quants use Gaussian Copulas to model the joint dependencies of asset returns. The determinant of the correlation matrix is essential for evaluating the joint density function of these assets.
\end{enumerate}

\subsection{B) Model Problem Setup: Multivariate Gaussian in ML}
Let us rigorously define the role of the determinant in Statistical Learning Theory. 
Define a state space $\mathbb{R}^n$ representing $n$ distinct features of a dataset (e.g., pixel intensities, financial indicators). Let $X$ be an $n$-dimensional random vector following a multivariate normal distribution, $X \sim \mathcal{N}(\mu, \Sigma)$, where $\mu \in \mathbb{R}^n$ is the mean vector and $\Sigma \in \mathbb{R}^{n \times n}$ is the covariance matrix.

The covariance matrix $\Sigma$ is symmetric and positive semi-definite. For the probability density function (PDF) to exist, $\Sigma$ must be strictly positive definite, which implies it is invertible and $\det(\Sigma) > 0$.

The PDF is given by:
\[ f(x) = \frac{1}{\sqrt{(2\pi)^n \det(\Sigma)}} \exp\left(-\frac{1}{2}(x-\mu)^T \Sigma^{-1} (x-\mu)\right) \]

\textbf{The Role of the Determinant:}
The term $(x-\mu)^T \Sigma^{-1} (x-\mu)$ is the squared Mahalanobis distance, measuring how far a point $x$ is from the mean $\mu$, scaled by the covariance. The contours of constant density are $n$-dimensional ellipsoids. 
The volume of these ellipsoids is directly proportional to $\sqrt{\det(\Sigma)}$. If $\det(\Sigma)$ approaches $0$, the volume collapses, meaning the data lies on a lower-dimensional hyperplane (perfect multicollinearity). The determinant ensures the total probability integrates to $1$ over $\mathbb{R}^n$ by acting as the precise geometric scaling factor for the volume of the feature space.

\newpage
\section{Part 7: Common Variations and Untested Concepts}

While the homework covered standard properties, a few advanced determinant concepts frequently appear on exams but were missing from this specific assignment.

\subsection{Missing Concept 1: Determinants of Block Matrices}
Often, you will encounter large matrices partitioned into smaller blocks. If a block matrix is upper or lower triangular in its blocks, its determinant is the product of the determinants of the diagonal blocks.
\textbf{Theorem:} If $M = \begin{bmatrix} A & B \\ 0 & C \end{bmatrix}$, where $A$ and $C$ are square matrices, then $\det(M) = \det(A)\det(C)$.

\textbf{Example:} Compute the determinant of $M = \begin{bmatrix} 2 & 5 & 9 & 1 \\ 0 & 3 & 8 & 4 \\ 0 & 0 & 4 & 2 \\ 0 & 0 & 1 & 3 \end{bmatrix}$.
\textbf{Solution:} We can partition this into $2 \times 2$ blocks:
$A = \begin{bmatrix} 2 & 5 \\ 0 & 3 \end{bmatrix}$, $B = \begin{bmatrix} 9 & 1 \\ 8 & 4 \end{bmatrix}$, $0 = \begin{bmatrix} 0 & 0 \\ 0 & 0 \end{bmatrix}$, $C = \begin{bmatrix} 4 & 2 \\ 1 & 3 \end{bmatrix}$.
Since the bottom-left block is all zeros, $\det(M) = \det(A)\det(C)$.
$\det(A) = (2)(3) - 0 = 6$.
$\det(C) = (4)(3) - (2)(1) = 12 - 2 = 10$.
$\det(M) = 6 \times 10 = 60$.

\subsection{Missing Concept 2: The Vandermonde Determinant}
A Vandermonde matrix is a specific type of matrix that appears in polynomial interpolation.
\textbf{Example:} $V = \begin{bmatrix} 1 & x_1 & x_1^2 \\ 1 & x_2 & x_2^2 \\ 1 & x_3 & x_3^2 \end{bmatrix}$.
The determinant of this matrix is always $\det(V) = (x_2 - x_1)(x_3 - x_1)(x_3 - x_2)$. This proves that as long as the points $x_1, x_2, x_3$ are distinct, the determinant is non-zero, meaning a unique interpolating polynomial always exists.

\newpage
\section{Part 8: Advanced Diagnostic Testing: "Find the Flaw"}
\textit{Instructions: You are the grader. Each of the following solutions contains a single, subtle conceptual or algebraic flaw. Identify the error, explain why it is wrong, and provide the correct solution.}

\subsection{Problem 1}
\textbf{Question:} Let $A = \begin{bmatrix} 1 & 0 \\ 0 & 1 \end{bmatrix}$ and $B = \begin{bmatrix} -1 & 0 \\ 0 & -1 \end{bmatrix}$. Find $\det(A + B)$.
\textbf{Flawed Solution:}
By the linearity property of determinants, we know that the determinant of a sum is the sum of the determinants.
$\det(A) = (1)(1) - 0 = 1$.
$\det(B) = (-1)(-1) - 0 = 1$.
Therefore, $\det(A + B) = \det(A) + \det(B) = 1 + 1 = 2$.

\textbf{Correction:}
\textbf{Flaw:} The determinant is \textit{not} a linear operator over matrix addition; $\det(A+B) \neq \det(A) + \det(B)$.
\textbf{Correct Solution:} First, compute the matrix sum $A+B = \begin{bmatrix} 1-1 & 0 \\ 0 & 1-1 \end{bmatrix} = \begin{bmatrix} 0 & 0 \\ 0 & 0 \end{bmatrix}$. The determinant of the zero matrix is $0$.

\subsection{Problem 2}
\textbf{Question:} Let $A$ be a $4 \times 4$ matrix with $\det(A) = 3$. Find $\det(2A)$.
\textbf{Flawed Solution:}
By the scalar multiplication property, we can factor the scalar out of the determinant.
$\det(2A) = 2 \det(A) = 2(3) = 6$.

\textbf{Correction:}
\textbf{Flaw:} When factoring a scalar out of an $n \times n$ matrix, it must be raised to the power of $n$, because the scalar multiplies every single row.
\textbf{Correct Solution:} $\det(2A) = 2^4 \det(A) = 16(3) = 48$.

\subsection{Problem 3}
\textbf{Question:} Evaluate $\begin{vmatrix} 0 & 2 \\ 3 & 4 \end{vmatrix}$ using row operations.
\textbf{Flawed Solution:}
We need a non-zero pivot in the top left. Let's swap Row 1 and Row 2.
$\begin{vmatrix} 0 & 2 \\ 3 & 4 \end{vmatrix} = \begin{vmatrix} 3 & 4 \\ 0 & 2 \end{vmatrix}$
The matrix is now upper triangular. The determinant is the product of the diagonal: $(3)(2) = 6$.

\textbf{Correction:}
\textbf{Flaw:} Swapping two rows multiplies the determinant by $-1$. The negative sign was omitted.
\textbf{Correct Solution:} $\begin{vmatrix} 0 & 2 \\ 3 & 4 \end{vmatrix} = - \begin{vmatrix} 3 & 4 \\ 0 & 2 \end{vmatrix} = - (3)(2) = -6$.

\subsection{Problem 4}
\textbf{Question:} Let $A = \begin{bmatrix} 1 & 2 \\ 3 & 4 \end{bmatrix}$. Apply the row operation $R_2 \leftarrow 2R_2 - 3R_1$ to create a zero in the bottom left, and find the determinant.
\textbf{Flawed Solution:}
Row replacement operations do not change the determinant.
$R_2 \leftarrow 2R_2 - 3R_1$ gives the new matrix $B = \begin{bmatrix} 1 & 2 \\ 2(3)-3(1) & 2(4)-3(2) \end{bmatrix} = \begin{bmatrix} 1 & 2 \\ 3 & 2 \end{bmatrix}$.
Wait, $2(3)-3(1) = 3$. Let's re-evaluate: $2(3) - 3(1) = 6 - 3 = 3$. This didn't create a zero.
Let's use $R_2 \leftarrow R_2 - 3R_1$.
$B = \begin{bmatrix} 1 & 2 \\ 0 & -2 \end{bmatrix}$.
Determinant is $(1)(-2) = -2$.

\textbf{Correction:}
\textbf{Flaw:} The flawed solution caught its own arithmetic error but missed the theoretical flaw in the prompt's suggested operation. The operation $R_2 \leftarrow 2R_2 - 3R_1$ is actually a combination of scaling ($R_2 \leftarrow 2R_2$) and replacement. Scaling changes the determinant!
\textbf{Correct Solution:} The standard replacement $R_2 \leftarrow R_2 - 3R_1$ preserves the determinant. $\det(A) = (1)(4) - (2)(3) = -2$. The matrix becomes $\begin{bmatrix} 1 & 2 \\ 0 & -2 \end{bmatrix}$, whose determinant is indeed $-2$.

\subsection{Problem 5}
\textbf{Question:} If $A$ is a square matrix such that $A^2 = 0$ (the zero matrix), prove that $A$ must be the zero matrix.
\textbf{Flawed Solution:}
We take the determinant of both sides: $\det(A^2) = \det(0)$.
Using the multiplicative property, $(\det(A))^2 = 0$.
Taking the square root, $\det(A) = 0$.
Since the determinant is zero, the matrix has no "volume" and all its entries must be zero. Therefore, $A = 0$.

\textbf{Correction:}
\textbf{Flaw:} A determinant of zero only implies the matrix is singular (not invertible), it does \textit{not} imply the matrix is the zero matrix.
\textbf{Correct Solution:} The premise is false. $A$ does not have to be the zero matrix. For example, let $A = \begin{bmatrix} 0 & 1 \\ 0 & 0 \end{bmatrix}$. $A \neq 0$, but $A^2 = \begin{bmatrix} 0 & 1 \\ 0 & 0 \end{bmatrix} \begin{bmatrix} 0 & 1 \\ 0 & 0 \end{bmatrix} = \begin{bmatrix} 0 & 0 \\ 0 & 0 \end{bmatrix}$. Such matrices are called nilpotent.

\end{document}